\section{Learning Convergence}

The training process of the Deep Reinforcement Learning agents was monitored to assess learning stability and convergence behavior, as well as to verify that improvements in the learned policies translated into reduced operational costs. Both the Soft Actor-Critic (SAC) and Proximal Policy Optimization (PPO) agents were trained for a total of 200{,}000 time steps. During training, the agents were periodically evaluated every 1{,}000 steps on a reserved validation dataset in order to track performance progression and mitigate overfitting.

The observed learning curves reveal distinct convergence characteristics for the two algorithms under the hyperparameter settings described in Section~5.1. The PPO agent exhibited a relatively smooth and monotonic improvement in average episodic reward throughout training. This behavior is consistent with the conservative update mechanism enforced by PPO’s clipped objective, which limits abrupt policy changes and promotes stable learning. The PPO agent reached a performance plateau after approximately 35{,}000 time steps, indicating convergence to a stable policy within the explored policy space.

In contrast, the SAC agent initially demonstrated higher variance in its episodic returns, reflecting the effect of entropy regularization, which encourages extensive exploration of the continuous action space during early training stages. As training progressed and the entropy temperature parameter ($\alpha$) adapted automatically, the SAC agent gradually reduced exploratory behavior and focused on higher-reward regions of the policy space. As a result, SAC converged to a higher final reward level compared to PPO.

The Robust SAC (R-SAC) agent, trained under domain randomization with injected forecast noise, exhibited a slower convergence rate due to the increased stochasticity of the training environment. Nevertheless, the agent ultimately stabilized at a performance level comparable to the standard SAC policy. This behavior indicates that the introduction of domain randomization increases training difficulty but does not prevent the agent from learning an effective control policy.

Overall, these results suggest that while PPO provides stable and predictable convergence, SAC’s entropy-regularized framework enables more extensive exploration and can yield improved long-term performance in continuous control settings. The slightly slower convergence of R-SAC reflects the expected trade-off between robustness and training efficiency.
